\documentclass[11pt,a4paper]{article}
\usepackage[left=2.2cm,right=2.2cm,top=1.8cm,bottom=1.8cm]{geometry}
\usepackage[T1]{fontenc}
\usepackage{lmodern}
\usepackage[english]{babel}
\usepackage[hidelinks]{hyperref}
\usepackage{parskip}
\pagestyle{empty}

\begin{document}

\begin{flushright}
Kostiantyn Pysanyi\\
Leipzig, Germany\\
\href{mailto:kostiantyn.pysanyi@gmail.com}{kostiantyn.pysanyi@gmail.com}\\
+49 1525 8447880
\end{flushright}

Wolt Enterprises Deutschland GmbH\\
Applied Science Hiring Team\\
Stralauer Allee 6\\
10245 Berlin

\begin{flushright}
4 September 2026
\end{flushright}

\textbf{Application for (Senior) Applied Scientist, Recommendations -- Job ID 3490497}

Dear Applied Science Hiring Team,

I am applying for the (Senior) Applied Scientist, Recommendations role because I enjoy the full loop of applied machine learning: understanding a difficult problem through its data, testing ideas carefully, and making the result reliable enough to use. Recommendations add a compelling scientific challenge because a useful result depends not only on model relevance, but also on context and what is available at that moment. I believe my experience developing, evaluating, and deploying ML systems for complex real-world data would bring a strong experimental and engineering foundation to this work.

At RAYLYTIC, I built an end-to-end 2D--3D registration pipeline for implant pose estimation, from mesh processing and alignment through validation and production deployment. The system improved the acceptance rate by 45\% and reduced error by 15\%. I also delivered production-grade segmentation, detection, and classification models, improving accuracy by 20--50\% over their baselines while reducing training time eightfold. To make this work reproducible, I implemented data versioning, experiment tracking, and automated evaluation. These projects required me to connect model decisions with clear metrics, investigate failures systematically, and turn research code into dependable services.

Earlier, as a Research Engineer, I designed data pipelines and training workflows for a ticket-routing classifier. The work transformed unstructured text into structured features and made dataset updates and model evaluation more efficient for industry partners. Across these projects, I learned to communicate technical choices clearly and to work across research, engineering, and domain boundaries.

Wolt particularly appeals to me because model decisions shape everyday customer choices while operating over changing local inventories. I would welcome the opportunity to contribute careful experimentation and production ML experience to recommendation systems used at the scale of Wolt and DoorDash.

I would be pleased to discuss how my background could contribute to Wolt's Recommendations team.

Sincerely,\\[0.8em]
Kostiantyn Pysanyi

\end{document}
